% rel_chapter_5_sec51-52.tex
% Relativistic Perturbation Equations with Rotation + Magnetic Field (Ch V, §51--52)
% Agent 21: rotation + magnetic field perturbation equations and stationary convection
% BDNK conversion: IS relaxation replaced by BDNK first-order
%   constitutive relations (Bemfica et al. 2018, Kovtun 2019)
%
% Conventions: (-,+,+,+) signature, c kept explicit, u^μ u_μ = -c²
% See RELATIVISTIC_CONVENTIONS.md for notation and macros.
% Macros assumed from rel_preamble.tex.

\chapter{Relativistic Rotation and Magnetic Field: Perturbation Equations
  and Stationary Convection}
\label{ch:rel-rotmag-pert}

In Chandrasekhar's Chapter~V, \S\S\,51--52, the combined effects of
rotation and a magnetic field on thermal convection are analysed in the
non-relativistic regime.  The perturbation equations that result
(equations~(21)--(43) of Chapter~V) couple the velocity, vorticity,
magnetic, and temperature perturbations through the Coriolis and
Lorentz forces simultaneously.  In the present chapter we derive the
fully relativistic counterparts of these equations, using the
BDNK (Bemfica--Disconzi--Noronha--Kovtun) first-order causal
dissipation formalism and covariant MHD developed in the relativistic
framework chapters.

Throughout, we keep the speed of light $c$ explicit, so that every
non-relativistic limit ($v/c\to 0$, $\vA/c\to 0$) can be read off
directly.  The gravitational field $\vect{g}$, the background
magnetic field $\vect{H}$, and the angular velocity
$\boldsymbol{\Omega}$ are taken to be parallel (both pointing in the
$z$-direction), as in the original treatment.

%======================================================================
\section{Relativistic perturbation equations (\S\,51)}
\label{sec:rel-51}
%======================================================================

\subsection{Background state and covariant rotating frame}
\label{sec:background-rot-mag}

We consider a layer of conducting fluid of depth~$d$, in uniform
rotation with angular velocity $\boldsymbol{\Omega} = \Omega\,\hat{z}$,
threaded by a uniform vertical magnetic field $\vect{H} = H\,\hat{z}$,
and heated from below with an adverse temperature gradient
$\beta = -\dd T_0/\dd z > 0$.

In the relativistic treatment, the equilibrium 4-velocity of the
rigidly rotating fluid is
\begin{equation}
  u^{\mu}_{(0)} = \Lf_\Omega\,(c,\;-\Omega y,\;\Omega x,\;0),
  \qquad
  \Lf_\Omega = \frac{1}{\sqrt{1 - \Omega^{2}r_{\perp}^{2}/c^{2}}},
  \label{eq:rel51-u0}
\end{equation}
where $r_\perp^2 = x^2+y^2$.  For $\Omega d \ll c$ (the regime of
interest), $\Lf_\Omega \simeq 1 + \order{\Omega^{2}d^{2}/c^{2}}$
and the corrections enter at second post-Newtonian order.

The background magnetic 4-vector in the fluid frame is
\begin{equation}
  b^{\mu}_{(0)} = \Lf_\Omega\,(0,\;0,\;0,\;B_0),
  \qquad
  b^{2}_{(0)} \equiv b^{\mu}_{(0)}\,b_{(0)\mu} = B_0^{2},
  \label{eq:rel51-b0}
\end{equation}
where $B_0 = \mu^{1/2} H/(4\pi)^{1/2}$ is the rest-frame field
strength in Gaussian units.

The total enthalpy density of the magnetised background fluid is
\begin{equation}
  \enthalpy^{*}_{0}
  = \frac{\varepsilon_{0} + p_{0}}{c^{2}} + \frac{b^{2}_{(0)}}{c^{2}}
  = \frac{\varepsilon_{0} + p_{0} + B_0^{2}}{c^{2}}.
  \label{eq:rel51-wstar}
\end{equation}

\subsection{Linearized relativistic MHD in the rotating frame}
\label{sec:lin-rmhd-rot}

We perturb all quantities about the equilibrium:
$u^{\mu} = u^{\mu}_{(0)} + \delta u^{\mu}$,
$b^{\mu} = b^{\mu}_{(0)} + \delta b^{\mu}$,
$\varepsilon = \varepsilon_{0} + \delta\varepsilon$,
$p = p_{0} + \delta p$, $T = T_0(z) + \Theta$.
Working in the corotating frame (where $u^{i}_{(0)}=0$ to leading
order), the spatial part of the perturbed velocity is
$\delta u^{\mu} = (0,\, u_1,\, u_2,\, w)$, and similarly for the
magnetic perturbation
$\delta b^{\mu} = (0,\, h_1,\, h_2,\, h_3)$.

The linearized covariant momentum equation, projected spatially with
$\Delta^{\mu\nu}$ and including Coriolis and Lorentz forces, gives
the \emph{relativistic} generalization of Chandrasekhar's
equation~(21):
\begin{multline}
  \enthalpy^{*}_{0}\,\frac{\partial u_{i}}{\partial t}
  = -\frac{\partial}{\partial x_{i}}\bigl(\delta p^{*}\bigr)
  + \rho_{0}\,g\,\alpha\,\Theta\,\delta_{i3}\,
    \relcorr{\vA^{2}/c^{2}}\\
  + \shearvisc\,\nabla^{2}u_{i}
  + 2\,\enthalpy^{*}_{0}\,\epsilon_{ijk}\,u_{j}\,\Omega_{k}\,
    \relcorr{\Omega^{2}d^{2}/c^{2}}\\
  + \frac{B_{0}}{c^{2}}\,
    \bigl(B_{0}\,\partial_{z}\,h_{i} - h_{3}\,\partial_{z}\,u_{i}^{(0)}\bigr)
  + \frac{B_{0}}{c^{2}}\,\partial_{z}\,h_{i},
  \label{eq:rel51-momentum}
\end{multline}
where $\delta p^{*} = \delta p + B_{0}\,h_{3}/c^{2}$ is the total
(fluid + magnetic) pressure perturbation, $\alpha$ is the coefficient
of thermal expansion, and we have used the BDNK shear
viscosity $\shearvisc$ in place of the Newtonian $\rho\nu$.  The
factors $\relcorr{\cdot}$ denote multiplicative relativistic
corrections.

Simplifying for the case $\Omega d \ll c$ and $\vA \ll c$, the
leading-order relativistic correction to the momentum equation is:
\begin{multline}
  \frac{\partial u_{i}}{\partial t}
  = -\frac{1}{\enthalpy^{*}_{0}}\,
    \frac{\partial(\delta p^{*})}{\partial x_{i}}
  + \frac{\rho_{0}\,g\,\alpha}{\enthalpy^{*}_{0}}\,\Theta\,\delta_{i3}
  + \frac{\shearvisc}{\enthalpy^{*}_{0}}\,\nabla^{2}u_{i}\\
  + 2\,\epsilon_{ijk}\,u_{j}\,\Omega_{k}
  + \frac{B_{0}}{\enthalpy^{*}_{0}\,c^{2}}\,
    B_{0}\,\frac{\partial h_{i}}{\partial z}.
  \label{eq:rel51-mom-simple}
\end{multline}
Note the replacement $\rho \to \enthalpy^{*}_{0}$ in the inertial
denominator: this is the hallmark of relativistic MHD, where the
magnetic field energy contributes to the effective inertia.

The linearized \emph{induction equation} in the corotating frame is
the relativistic counterpart of equation~(22):
\begin{equation}
  \frac{\partial h_{i}}{\partial t}
  = B_{0}\,\frac{\partial u_{i}}{\partial z}
  + \eta\,\nabla^{2}h_{i}
  + \order{\vA^{2}/c^{2}}\text{ corrections},
  \label{eq:rel51-induction}
\end{equation}
where $\eta$ is the magnetic diffusivity.  The relativistic
correction arises from the coupling between the magnetic 4-vector
evolution and the 4-acceleration; it is suppressed by $\vA^{2}/c^{2}$.

The linearized \emph{energy equation} with BDNK causal heat
transport replaces the classical equation~(23):
\begin{equation}
  \frac{\partial\Theta}{\partial t}
  = \beta\,w + \kappa\,\nabla^{2}\Theta
  + \order{v^{2}/c^{2}},
  \label{eq:rel51-energy}
\end{equation}
In the BDNK formalism \citep{Bemfica2018,Kovtun2019}, the heat
equation retains its standard first-order form---no relaxation time
$\tauq$ or telegrapher modification appears.  Causality of thermal
signal propagation is guaranteed by the BDNK frame coefficients,
which ensure strong hyperbolicity of the full conservation-law system
with all characteristic speeds bounded by~$c$.

The \emph{divergence-free conditions} remain:
\begin{equation}
  \frac{\partial u_{i}}{\partial x_{i}} = 0
  \quad\text{and}\quad
  \frac{\partial h_{i}}{\partial x_{i}} = 0.
  \label{eq:rel51-divfree}
\end{equation}

\subsection{Vorticity and double-curl equations}
\label{sec:vort-doublecurl}

Following the same procedure as in the classical theory---taking the
curl and double curl of the momentum equation and extracting the
$z$-components---we obtain the relativistic generalizations of
equations~(28)--(29).

Define the effective kinematic viscosity and Alfv\'en speed:
\begin{equation}
  \nu_{\mathrm{eff}}
  = \frac{\shearvisc}{\enthalpy^{*}_{0}},
  \qquad
  \vA^{2}
  = \frac{B_{0}^{2}\,c^{2}}{4\pi(\varepsilon_{0}+p_{0})+B_{0}^{2}}.
  \label{eq:rel51-nueff-vA}
\end{equation}
Note that $\vA < c$ is guaranteed by the dominant energy condition.

The $z$-component of the vorticity equation is:
\begin{equation}
  \frac{\partial\zeta}{\partial t}
  = \nu_{\mathrm{eff}}\,\nabla^{2}\zeta
  + 2\Omega\,\frac{\partial w}{\partial z}
  + \frac{\vA^{2}}{c^{2}}\,
    \frac{B_{0}}{4\pi\,\enthalpy^{*}_{0}}\,
    \frac{\partial\xi}{\partial z},
  \label{eq:rel51-zeta}
\end{equation}
where $\zeta$ and $\xi$ are the $z$-components of the vorticity and
current perturbation, respectively.

The double-curl ($z$-component of $\nabla^{2}\vect{u}$) equation:
\begin{multline}
  \frac{\partial}{\partial t}(\nabla^{2}w)
  = \frac{\rho_{0}\,g\,\alpha}{\enthalpy^{*}_{0}}
    \left(\frac{\partial^{2}\Theta}{\partial x^{2}}
    +\frac{\partial^{2}\Theta}{\partial y^{2}}\right)
  + \nu_{\mathrm{eff}}\,\nabla^{4}w\\
  - 2\Omega\,\frac{\partial\zeta}{\partial z}
  + \frac{\vA^{2}}{c^{2}}\,
    \frac{B_{0}}{4\pi\,\enthalpy^{*}_{0}}\,
    \frac{\partial}{\partial z}\,\nabla^{2}h_{3}.
  \label{eq:rel51-doublecurl}
\end{multline}

\subsection{Normal mode decomposition}
\label{sec:normal-modes}

Analysing the perturbations into normal modes
$\sim \exp[\mathrm{i}(k_{x}x + k_{y}y) + \sigma t/d^{2}]$
with $a^{2} = k_{x}^{2}+k_{y}^{2}$ (in units of $1/d$), and
introducing the non-dimensional vertical coordinate $z' = z/d$
with $D = \dd/\dd z'$, the system
\eqref{eq:rel51-energy}--\eqref{eq:rel51-doublecurl} becomes:

\begin{equation}
  \bigl(D^{2}-a^{2}-p_{1}\sigma\bigr)\,\Theta
  = -\frac{\beta\,d^{2}}{\kappa}\,W,
  \label{eq:rel51-35}
\end{equation}
\begin{equation}
  (D^{2}-a^{2}-p_{2}\,\sigma)\,K
  = -\frac{B_{0}\,d}{\eta}\,DW,
  \label{eq:rel51-36}
\end{equation}
\begin{equation}
  (D^{2}-a^{2}-p_{2}\,\sigma)\,X
  = -\frac{B_{0}\,d}{\eta}\,DZ,
  \label{eq:rel51-37}
\end{equation}
\begin{equation}
  (D^{2}-a^{2}-\sigma)\,Z
  = -\frac{2\Omega\,d^{2}}{\nu_{\mathrm{eff}}}\,DW
  - \frac{\vA^{2}\,d}{c^{2}\,\nu_{\mathrm{eff}}}\,
    \frac{B_{0}}{4\pi\,\enthalpy^{*}_{0}}\,DX,
  \label{eq:rel51-38}
\end{equation}
and
\begin{multline}
  (D^{2}-a^{2})(D^{2}-a^{2}-\sigma)\,W
  + \frac{\vA^{2}\,d}{c^{2}\,\nu_{\mathrm{eff}}}\,
    \frac{B_{0}}{4\pi\,\enthalpy^{*}_{0}}\,
    D(D^{2}-a^{2})\,K\\
  - \frac{2\Omega\,d^{2}}{\nu_{\mathrm{eff}}}\,DZ
  = \frac{\rho_{0}\,g\,\alpha\,d^{2}}
         {\enthalpy^{*}_{0}\,\nu_{\mathrm{eff}}}\,a^{2}\,\Theta.
  \label{eq:rel51-39}
\end{multline}

\subsection{Relativistic dimensionless parameters}
\label{sec:rel-dimless}

We now introduce the relativistic dimensionless parameters.  The
\emph{relativistic Rayleigh number} is
\begin{equation}
  \Rarel
  = \frac{\rho_{0}\,g\,\alpha\,\beta\,d^{4}}
         {\enthalpy^{*}_{0}\,\nu_{\mathrm{eff}}\,\kappa}
  = \mathrm{Ra}\times
    \frac{\rho_{0}}{\enthalpy^{*}_{0}}
    \times\frac{\nu}{\nu_{\mathrm{eff}}}.
  \label{eq:rel51-Rarel}
\end{equation}
The \emph{relativistic Taylor number} is
\begin{equation}
  \Tarel
  = \frac{4\Omega^{2}\,d^{4}}{\nu_{\mathrm{eff}}^{2}}
  = T\times\left(\frac{\nu}{\nu_{\mathrm{eff}}}\right)^{\!2},
  \label{eq:rel51-Tarel}
\end{equation}
and the \emph{relativistic Chandrasekhar number} is
\begin{equation}
  \Qrel
  = \frac{B_{0}^{2}\,d^{2}}
         {4\pi\,\enthalpy^{*}_{0}\,\nu_{\mathrm{eff}}\,\eta}
  = Q\times
    \frac{\rho_{0}}{\enthalpy^{*}_{0}}
    \times\frac{\nu}{\nu_{\mathrm{eff}}}.
  \label{eq:rel51-Qrel}
\end{equation}
In each case, the classical parameter is multiplied by a factor
$\rho_{0}/\enthalpy^{*}_{0} < 1$, reflecting the enhanced inertia
from both thermal energy ($\varepsilon_0 + p_0 > \rho_0 c^2$) and
magnetic field energy.  In the non-relativistic limit
$\enthalpy^{*}_{0} \to \rho_{0}$ and $\nu_{\mathrm{eff}}\to\nu$,
so all three reduce to the classical values.

\subsection{Elimination and the master equation}
\label{sec:master-eqn}

Eliminating $K$, $X$, $Z$, and $\Theta$ between
equations~\eqref{eq:rel51-35}--\eqref{eq:rel51-39} by the same
algebraic procedure as in the classical theory
(cf.\ equations~(40)--(43) of Chapter~V), we obtain the
\emph{relativistic master equation}:
\begin{multline}
  \bigl(D^{2}-a^{2}-p_{1}\sigma\bigr)\\
  \times\bigl\{(D^{2}-a^{2})
    \bigl[(D^{2}-a^{2}-\sigma)(D^{2}-a^{2}-p_{2}\sigma)
      -\Qrel\,D^{2}\bigr]^{2}\\
  +\;\Tarel\,D^{2}(D^{2}-a^{2}-p_{2}\sigma)^{2}\bigr\}\,W\\
  = -\Rarel\,a^{2}
    \bigl[(D^{2}-a^{2}-\sigma)(D^{2}-a^{2}-p_{2}\sigma)
      -\Qrel\,D^{2}\bigr]\\
  \quad\times(D^{2}-a^{2}-p_{2}\sigma)\,W.
  \label{eq:rel51-master}
\end{multline}
This is formally identical to Chandrasekhar's equation~(43), with the
replacements $R\to\Rarel$, $T\to\Tarel$, $Q\to\Qrel$.  In the BDNK
formalism, no additional Israel--Stewart relaxation term appears in the
thermal operator---causality is ensured by the frame coefficients, not
by modifying the PDE structure.

\begin{causalitycheck}
\textbf{Thermal modes:} In the BDNK formalism, the heat operator
remains first-order in~$\sigma$: $(D^{2}-a^{2}-p_{1}\sigma)$.
Causality of thermal signal propagation is guaranteed by the BDNK
frame coefficients, which ensure strong hyperbolicity of the full
conservation-law system with all characteristic speeds $\le c$
\citep{Bemfica2018,Kovtun2019}.

\textbf{Alfv\'en and magnetosonic modes:} The magnetic sector
produces the relativistic Alfv\'en speed
$\vA^{2} = B_{0}^{2}c^{2}/[4\pi(\varepsilon_{0}+p_{0})+B_{0}^{2}]
< c^{2}$,
and fast magnetosonic speed
$\vf^{2} = \cs^{2}+\vA^{2}-\cs^{2}\vA^{2}/c^{2} < c^{2}$.

\textbf{Inertial (Coriolis) modes:} Rotation introduces inertial
waves with frequency $\omega \le 2\Omega$.  For $\Omega d \ll c$
(always satisfied for laboratory and astrophysical convection layers),
the group speed $v_{g} \sim \Omega d \ll c$.

\textbf{Conclusion:} All perturbation modes propagate at speeds
$\le c$.  The system~\eqref{eq:rel51-master} is causal.
\end{causalitycheck}


%======================================================================
\section{Stationary convection --- relativistic solutions (\S\,52)}
\label{sec:rel-52}
%======================================================================

\subsection{Marginal state: $\sigma = 0$}

When instability sets in as stationary (non-oscillatory) convection,
the marginal state has $\sigma = 0$.  Setting $\sigma=0$ in the
normal-mode equations~\eqref{eq:rel51-35}--\eqref{eq:rel51-39} gives
the relativistic counterparts of equations~(44)--(48):
\begin{equation}
  (D^{2}-a^{2})\,\Theta
  = -\frac{\beta\,d^{2}}{\kappa}\,W,
  \label{eq:rel52-44}
\end{equation}
\begin{equation}
  \bigl[(D^{2}-a^{2})^{2}-\Qrel\,D^{2}\bigr]Z
  = -\frac{2\Omega\,d^{2}}{\nu_{\mathrm{eff}}}\,
    D(D^{2}-a^{2})\,W,
  \label{eq:rel52-45}
\end{equation}
\begin{equation}
  (D^{2}-a^{2})\,K
  = -\frac{B_{0}\,d}{\eta}\,DW,
  \label{eq:rel52-46}
\end{equation}
\begin{equation}
  (D^{2}-a^{2})\,X
  = -\frac{B_{0}\,d}{\eta}\,DZ,
  \label{eq:rel52-47}
\end{equation}
and the combined equation:
\begin{multline}
  (D^{2}-a^{2})
  \bigl\{\bigl[(D^{2}-a^{2})^{2}-\Qrel\,D^{2}\bigr]^{2}
  +\Tarel\,D^{2}(D^{2}-a^{2})^{2}\bigr\}\,W\\
  = -\Rarel\,a^{2}\,
    \bigl[(D^{2}-a^{2})^{2}-\Qrel\,D^{2}\bigr]\,W.
  \label{eq:rel52-48}
\end{multline}

The boundary conditions are identical in form to the classical ones
(equations~(49)--(51) of Chapter~V), since they express kinematic,
thermal, and electromagnetic matching conditions that are
frame-independent at leading order.

Note that at $\sigma=0$ the dissipative terms vanish regardless of
the causal formalism used (BDNK or Israel--Stewart): the
marginal-stability criterion for stationary convection is
\emph{insensitive} to the specific causal modification of heat
transport.  This is exactly as expected---dissipative corrections
affect only the time-dependent behaviour.

\subsection{Two free boundaries: relativistic critical Rayleigh number}
\label{sec:two-free-bdry}

For two stress-free, perfectly permeable boundaries ($W = D^{2}W = 0$,
$DZ = 0$ at $z=0,1$), the proper solutions are sinusoidal, exactly
as in the classical case.  Substituting $W = W_{0}\sin(n\pi z)$ into
equation~\eqref{eq:rel52-48}, we find the relativistic characteristic
equation:
\begin{equation}
  \boxed{
  R_{\mathrm{rel}}
  = \frac{(n^{2}\pi^{2}+a^{2})\,
    \bigl\{\bigl[(n^{2}\pi^{2}+a^{2})^{2}
      +\Qrel\,n^{2}\pi^{2}\bigr]^{2}
    +\Tarel\,n^{2}\pi^{2}(n^{2}\pi^{2}+a^{2})^{2}\bigr\}}
    {a^{2}\,\bigl[(n^{2}\pi^{2}+a^{2})^{2}
      +\Qrel\,n^{2}\pi^{2}\bigr]}.
  }
  \label{eq:rel52-57}
\end{equation}
This has the \emph{same functional form} as Chandrasekhar's
equation~(57), with the substitutions $R\to R_{\mathrm{rel}}$,
$Q\to\Qrel$, $T\to\Tarel$.  The instability first sets in for
$n=1$.

Introducing $x = a^{2}/\pi^{2}$,
$Q_{1,\mathrm{rel}} = \Qrel/\pi^{2}$, and
$T_{1,\mathrm{rel}} = \Tarel/\pi^{4}$, we write:
\begin{equation}
  R_{\mathrm{rel}} = \pi^{4}\,\frac{(1+x)\,
    \bigl\{\bigl[(1+x)^{2}+Q_{1,\mathrm{rel}}\bigr]^{2}
    +T_{1,\mathrm{rel}}(1+x)\bigr\}}
    {x\,\bigl[(1+x)^{2}+Q_{1,\mathrm{rel}}\bigr]}.
  \label{eq:rel52-59}
\end{equation}
The critical Rayleigh number is obtained by minimising over~$x$:
\begin{equation}
  2x^{3}+3x^{2}-1
  = Q_{1,\mathrm{rel}}
  + T_{1,\mathrm{rel}}\,
    \frac{(1+x)^{4}-Q_{1,\mathrm{rel}}(x^{2}-1)}
         {\bigl[(1+x)^{2}+Q_{1,\mathrm{rel}}\bigr]^{2}}.
  \label{eq:rel52-61}
\end{equation}

\begin{relcorrection}
The critical Rayleigh number for stationary convection with both
rotation and magnetic field in the relativistic theory is related to
the classical value by
\begin{equation}
  R_{c,\mathrm{rel}}
  = R_{c,\mathrm{classical}}\;\times\;
  \frac{\enthalpy^{*}_{0}}{\rho_{0}}\;\times\;
  \frac{\nu_{\mathrm{eff}}}{\nu}
  = R_{c,\mathrm{classical}}\;\times\;
  \left(1 + \frac{p_{0}+\varepsilon_{\mathrm{int}}+B_{0}^{2}}
                  {\rho_{0}\,c^{2}}\right)
  \left(1 + \order{\vA^{2}/c^{2}}\right).
  \label{eq:rel52-Rc-relation}
\end{equation}
The relativistic correction \emph{raises} the critical Rayleigh number:
it is harder to destabilise the layer when the effective inertia is
enhanced by internal energy and magnetic field energy.  This effect is
small for laboratory fluids ($p_{0}/\rho_{0}c^{2}\sim 10^{-10}$) but
can be significant in neutron star envelopes and accretion discs where
$p_{0}/\rho_{0}c^{2}\sim 0.1$--$1$.
\end{relcorrection}

\subsection{Competition and cooperation of rotation and magnetic field}
\label{sec:compete-cooperate}

The structure of equation~\eqref{eq:rel52-59} reveals the same
interplay between rotation and magnetic field as in the classical
theory, but with quantitative modifications:

\paragraph{Magnetic field alone ($\Tarel = 0$).}
The critical Rayleigh number increases monotonically with $\Qrel$, and
for large $\Qrel$,
\begin{equation}
  R_{c,\mathrm{rel}} \sim \pi^{2}\,\Qrel
  \qquad (\Tarel = 0,\;\Qrel\gg 1).
  \label{eq:rel52-Qalone}
\end{equation}

\paragraph{Rotation alone ($\Qrel = 0$).}
The critical Rayleigh number increases with $\Tarel$, and for large
$\Tarel$,
\begin{equation}
  R_{c,\mathrm{rel}}
  \sim \left(\frac{27\pi^{4}}{4}\right)^{1/3}\!\Tarel^{2/3}
  \qquad (\Qrel = 0,\;\Tarel\gg 1).
  \label{eq:rel52-Talone}
\end{equation}

\paragraph{Both acting ($\Qrel > 0$, $\Tarel > 0$).}
The central result---which persists unchanged in the relativistic
regime---is that a magnetic field can \emph{partially relieve} the
rotational inhibition of convection.  For fixed $\Tarel$, increasing
$\Qrel$ from zero initially \emph{decreases} $R_{c,\mathrm{rel}}$
before eventually increasing it.  The physical mechanism is unchanged:
the magnetic field couples the Coriolis-deflected horizontal motions
back to the vertical, effectively short-circuiting the Taylor--Proudman
constraint.

The novel relativistic feature is that the competition is modulated by
the ratio $\vA^{2}/c^{2}$:
\begin{itemize}
\item When $\vA^{2}/c^{2}$ is non-negligible, the magnetic field
  contributes to the effective inertia (via $\enthalpy^{*}_{0}$),
  which \emph{simultaneously reduces} both $\Qrel$ and $\Tarel$
  relative to their classical values.  The net effect depends on
  which reduction dominates.
\item For $B_0^{2} \sim 4\pi(\varepsilon_{0}+p_{0})$ (magnetically
  dominated regime), the effective inertia is roughly doubled, and
  both $\Qrel$ and $\Tarel$ are reduced by a factor of~2.
  Since $R_c$ depends on $\Tarel^{2/3}$ but linearly on $\Qrel$ in
  the respective asymptotic limits, the rotational stabilization is
  reduced faster than the magnetic stabilization.
\end{itemize}

\subsection{Two-minimum phenomenon in the relativistic regime}
\label{sec:two-minima}

The remarkable two-minimum structure of the marginal curve
$R_{\mathrm{rel}}(a)$ discovered by Chandrasekhar persists in the
relativistic theory, since the functional form of
equation~\eqref{eq:rel52-59} is preserved.  For certain ranges of
$Q_{1,\mathrm{rel}}$ and $T_{1,\mathrm{rel}}$, the marginal curve
has two local minima: one at large~$a$ (elongated cells) and one at
small~$a$ (flattened cells).

The \emph{transition value} of $Q_{1,\mathrm{rel}}$ at which the
global minimum jumps discontinuously from one branch to the other is
shifted relative to the classical value.  Specifically, for a given
$T_{1,\mathrm{rel}}$, the transition occurs at
\begin{equation}
  Q_{1,\mathrm{rel}}^{(\mathrm{trans})}
  = Q_{1,\mathrm{classical}}^{(\mathrm{trans})}
    \times \frac{\rho_{0}}{\enthalpy^{*}_{0}}
    \times \frac{\nu}{\nu_{\mathrm{eff}}}.
  \label{eq:rel52-Qtrans}
\end{equation}
Since $\rho_{0}/\enthalpy^{*}_{0} < 1$, the transition to flattened
cells occurs at a \emph{lower} magnetic field strength in the
relativistic regime.

\subsection{Modified stability boundaries}
\label{sec:modified-boundaries}

Collecting results, the relativistic stability boundary for stationary
convection with rotation and magnetic field is:
\begin{equation}
  R_{c,\mathrm{rel}} = R_{c,\mathrm{rel}}(\Tarel, \Qrel),
  \label{eq:rel52-boundary}
\end{equation}
with the functional form given by minimising
equation~\eqref{eq:rel52-59} over~$x$.  The boundary surface in the
$(\Rarel, \Tarel, \Qrel)$ parameter space has the same topology as the
classical surface, but is shifted by the relativistic correction
factor.

\begin{relcorrection}
\textbf{Summary of relativistic modifications to stability
boundaries:}
\begin{enumerate}
\item The effective inertia $\enthalpy^{*}_{0}$ replaces $\rho_{0}$,
  raising $R_{c}$ and lowering $Q$ and~$T$ relative to classical
  values.
\item Dissipative (BDNK) corrections do not affect the
  stationary-convection boundary ($\sigma=0$).
\item The Alfv\'en speed is bounded: $\vA < c$, preventing unphysical
  divergence of the Lorentz force coupling at strong fields.
\item The two-minimum phenomenon persists but the transition field
  strength is reduced by $\rho_{0}/\enthalpy^{*}_{0}$.
\item All perturbation mode speeds remain subluminal.
\end{enumerate}
\end{relcorrection}

\subsection{Non-relativistic limit}
\label{sec:NR-limit}

As a consistency check, we verify that all results reduce to
Chandrasekhar's in the non-relativistic limit.  Setting
$\varepsilon_{0}\to\rho_{0}c^{2}$, $p_{0}\ll\rho_{0}c^{2}$,
$B_{0}^{2}\ll 4\pi\rho_{0}c^{2}$ (and the BDNK frame corrections
vanish at $\order{v^{2}/c^{2}}$):
\begin{align}
  \enthalpy^{*}_{0} &\NRlimit \rho_{0},\\
  \nu_{\mathrm{eff}} &\NRlimit \nu,\\
  \Rarel &\NRlimit \mathrm{Ra} = \frac{g\alpha\beta d^{4}}{\nu\kappa},\\
  \Tarel &\NRlimit T = \frac{4\Omega^{2}d^{4}}{\nu^{2}},\\
  \Qrel &\NRlimit Q = \frac{\mu H^{2}d^{2}}{4\pi\rho\nu\eta},\\
  \vA^{2} &\NRlimit \frac{B^{2}}{4\pi\rho},
\end{align}
and equation~\eqref{eq:rel52-57} reduces identically to
Chandrasekhar's equation~(57).  The master equation
\eqref{eq:rel51-master} reduces to equation~(43).\qed

\begin{causalitycheck}
\textbf{Final causality summary for all perturbation modes in
\S\S\,51--52:}
\begin{enumerate}
\item \textbf{Thermal diffusion:} In BDNK, the first-order heat
  equation has causal signal propagation guaranteed by the frame
  coefficients; all thermal characteristic speeds $\le c$
  \citep{Bemfica2018,Kovtun2019}.
\item \textbf{Alfv\'en waves:} Propagation speed
  $\vA = B_{0}c/\sqrt{4\pi(\varepsilon_{0}+p_{0})+B_{0}^{2}} < c$.
\item \textbf{Fast magnetosonic:} $\vf^{2}=\cs^{2}+\vA^{2}
  -\cs^{2}\vA^{2}/c^{2} < c^{2}$.
\item \textbf{Slow magnetosonic:} $\vs^{2}=\cs^{2}\vA^{2}/\vf^{2}
  < \min(\cs^{2},\vA^{2}) < c^{2}$.
\item \textbf{Inertial waves:} Group speed $\sim\Omega d\ll c$.
\item \textbf{Viscous diffusion:} In the BDNK formalism, the viscous
  signal speed is bounded by~$c$ through the coupled inequalities on
  $\shearvisc$ and the frame coefficients
  \citep{BemficaDisconziNoronhaKovtun2023}; no separate relaxation
  time $\taupi$ is required.
\end{enumerate}
All six mode families satisfy the causality bound $v\le c$.
\end{causalitycheck}
